\documentclass{article}
% ---------------------------------------------------------------
% Póster de revisión de literatura (A0 vertical)
% Finanzas I (055D), UNCP, 2026-II. Compilar con pdfLaTeX.
% ---------------------------------------------------------------
\usepackage[paperwidth=841mm,paperheight=1189mm,margin=0mm]{geometry}
\usepackage[utf8]{inputenc}
\usepackage[T1]{fontenc}
\IfFileExists{spanish.ldf}{\usepackage[spanish,es-noshorthands,es-nodecimaldot]{babel}}{}
\usepackage{helvet}
\renewcommand{\familydefault}{\sfdefault}
\usepackage[letterspace=120]{microtype}
\usepackage{anyfontsize}
\usepackage{graphicx}
\usepackage{xcolor}
\usepackage[most]{tcolorbox}
\usepackage{enumitem}
\usepackage{multicol}
\usepackage{xurl}
\usepackage[hidelinks]{hyperref}
\usepackage{qrcode}
\urlstyle{same}
\pagestyle{empty}
\setlength{\parindent}{0pt}
\hyphenpenalty=10000 \exhyphenpenalty=10000

% ---------------- Paleta ----------------
\definecolor{navy}{HTML}{8E0C11}
\definecolor{navyLight}{HTML}{B8121A}
\definecolor{gold}{HTML}{00A3E0}
\definecolor{acento}{HTML}{A6E1F7}
\definecolor{goldSoft}{HTML}{E8F6FC}
\definecolor{teal}{HTML}{353435}
\definecolor{tealSoft}{HTML}{F1EEEE}
\definecolor{ink}{HTML}{231F20}
\definecolor{muted}{HTML}{5E5757}
\definecolor{line}{HTML}{E5D7D7}
\definecolor{page}{HTML}{F6F1F1}
\definecolor{mist}{HTML}{F7D3D5}
\pagecolor{page}

% ---------------- Tamaños (ajustados por póster) ----------------
\newcommand{\bodysize}{\fontsize{31}{40.9}\selectfont}
\newcommand{\cardtitlesize}{\fontsize{36.6}{43.9}\selectfont}
\newcommand{\smallsize}{\fontsize{25.4}{33.0}\selectfont}

% ---------------- Estilos de tarjeta ----------------
\tcbset{
  card/.style={enhanced, sharp corners, colback=white, colframe=line, boxrule=0.9mm,
    left=14mm, right=14mm, top=12mm, bottom=12mm, fontupper=\bodysize\color{ink}\raggedright,
    overlay={\fill[#1] (frame.north west) rectangle ([yshift=-5mm]frame.north east);}},
  card/.default=navy,
}
\newcommand{\cardtitle}[2][navy]{{\color{#1}\cardtitlesize\bfseries\textls[110]{\MakeUppercase{#2}}}\par\vspace{7mm}}
\newcommand{\bul}{\raisebox{0.18em}{\color{gold}\rule{0.5em}{0.5em}}}
\setlist[itemize]{label=\bul, leftmargin=1.25em, labelsep=0.6em, itemsep=0.55em, topsep=0pt, parsep=0pt}
\newcommand{\pill}[3]{\tcbox[on line, enhanced, arc=7mm, boxrule=0.8mm, colback=#1, colframe=#2,
  left=7mm, right=7mm, top=3mm, bottom=3mm, fontupper=\fontsize{26}{30}\selectfont\bfseries]{#3}}

\begin{document}

% ======================= ENCABEZADO =======================
\begin{tcolorbox}[enhanced, sharp corners, colback=navy, boxrule=0pt, width=\paperwidth,
  left=34mm, right=34mm, top=20mm, bottom=20mm,
  overlay={\fill[gold] (frame.south west) rectangle ([yshift=5mm]frame.south east);
           \fill[navyLight] ([xshift=-160mm]frame.north east) -- (frame.north east) -- ([yshift=-160mm]frame.north east) -- cycle;}]
\begin{minipage}[c]{92mm}\centering\includegraphics[width=\linewidth]{Escudo_UNCP.png}\end{minipage}\hfill
\begin{minipage}[c]{555mm}\centering
{\color{acento}\fontsize{25}{30}\selectfont\bfseries\textls[160]{REVISIÓN DE LITERATURA \quad|\quad FINANZAS I \quad|\quad TEMA 61}\par}
\vspace{10mm}
{\color{white}\fontsize{66}{74.6}\selectfont\bfseries What drives stock returns across countries? Insights from machine learning models\par}
\vspace{8mm}
{\color{mist}\fontsize{38}{46.4}\selectfont\itshape ¿Qué impulsa los retornos bursátiles entre países? Aportes de los modelos de aprendizaje automático\par}
\vspace{10mm}
{\color{white}\fontsize{29}{35}\selectfont\bfseries Nusret Cakici y Adam Zaremba\par}
\vspace{9mm}
\pill{white}{white}{\color{navy}Q1}\hspace{6mm}\pill{navy}{white}{\color{white}\mdseries \textit{International Review of Financial Analysis}, 96(Parte A), 103569 (2024)}
\end{minipage}\hfill
\begin{minipage}[c]{92mm}\centering\includegraphics[width=\linewidth]{Logo_FE.png}\end{minipage}
\end{tcolorbox}
\vspace{-1mm}
\begin{tcolorbox}[enhanced, sharp corners, colback=goldSoft, boxrule=0pt, width=\paperwidth,
  left=34mm, right=34mm, top=8mm, bottom=8mm, halign=center]
{\color{navy}\fontsize{27}{36}\selectfont \textbf{Póster elaborado por:} Meglinho Sneyder Rojas Poma \enspace\textbar\enspace Código 2024200522B\par \fontsize{24}{32}\selectfont Escuela Profesional de Economía, Universidad Nacional del Centro del Perú \enspace\textbar\enspace Finanzas I (055D), 2026-II\par}
\end{tcolorbox}

\vfill
% ======================= CUERPO =======================
\hspace*{34mm}\begin{minipage}{773mm}
\begin{tcolorbox}[enhanced, sharp corners, colback=navy, boxrule=0pt, left=16mm, right=16mm, top=12mm, bottom=11mm,
  fontupper=\bodysize\raggedright, overlay={\fill[gold] (frame.north west) rectangle ([yshift=-5mm]frame.north east);}]
\begin{minipage}[t]{174.8mm}\raggedright{\color{acento}\smallsize\bfseries\textls[110]{\MakeUppercase{Mercado}}}\par\vspace{2mm}{\color{white}71 mercados bursátiles}\end{minipage}\hfill\begin{minipage}[t]{174.8mm}\raggedright{\color{acento}\smallsize\bfseries\textls[110]{\MakeUppercase{Muestra}}}\par\vspace{2mm}{\color{white}87 predictores}\end{minipage}\hfill\begin{minipage}[t]{174.8mm}\raggedright{\color{acento}\smallsize\bfseries\textls[110]{\MakeUppercase{Periodo}}}\par\vspace{2mm}{\color{white}1985 a 2021}\end{minipage}\hfill\begin{minipage}[t]{174.8mm}\raggedright{\color{acento}\smallsize\bfseries\textls[110]{\MakeUppercase{Método}}}\par\vspace{2mm}{\color{white}Aprendizaje automático}\end{minipage}
\end{tcolorbox}
\end{minipage}

\vfill
\hspace*{34mm}\begin{minipage}{773mm}
\begin{tcbraster}[raster columns=3, raster equal height, raster column skip=16mm]
\begin{tcolorbox}[card]\cardtitle{Contexto}La lista de variables que predicen los retornos de los países crece cada año: valor, momentum, beta, riesgo político o tasas de interés. Casi siempre se estudian por separado.\end{tcolorbox}
\begin{tcolorbox}[card]\cardtitle{Problema}No se sabía qué variables pesan más cuando se analizan todas juntas, ni si la predictibilidad es mayor en mercados pequeños y poco integrados.\end{tcolorbox}
\begin{tcolorbox}[card]\cardtitle{Objetivo}Usar aprendizaje automático para explicar las diferencias de retorno entre los mercados bursátiles de distintos países.\end{tcolorbox}
\end{tcbraster}
\end{minipage}

\vfill
\hspace*{34mm}\begin{minipage}{773mm}
\begin{tcolorbox}[card]\cardtitle{Metodología}
\begin{minipage}[t]{0.48\linewidth}
\begin{itemize}
\item Se comparan MCO, reducción de dimensión, regresiones regularizadas, máquinas de vectores de soporte, árboles y redes neuronales.
\item Se combinan los pronósticos.
\item Se mide la importancia de cada variable en distintos horizontes.
\end{itemize}
\end{minipage}\hfill
\begin{minipage}[t]{0.48\linewidth}
\begin{itemize}
\item Los países se ordenan en quintiles según el retorno pronosticado y se compra el superior y se vende el inferior.
\item La capacidad predictiva fuera de muestra se mide con el R$^2$ fuera de muestra.
\end{itemize}
\end{minipage}
\end{tcolorbox}
\end{minipage}

\vfill
\hspace*{34mm}\begin{minipage}{773mm}
\begin{tcolorbox}[card]\cardtitle{Figura clave del estudio}
\begin{minipage}[c]{420.0mm}\centering
\fcolorbox{line}{white}{\includegraphics[width=\dimexpr420.0mm-2\fboxsep-2\fboxrule\relax]{figura.png}}
\end{minipage}\hfill
\begin{minipage}[c]{\dimexpr\linewidth-420.0mm-18mm\relax}\raggedright
{\color{navy}\bfseries Figura 8.}\enspace{\itshape Retornos acumulados de carteras largas y cortas formadas con los pronósticos de diez modelos de aprendizaje automático (OLS, PCA, PLS, LASSO, ENET, SVM, RF, GBRT, FFNN y COMB).}\par\vspace{8mm}
{\color{teal}\smallsize\bfseries\textls[110]{CÓMO LEERLA}}\par\vspace{4mm}
\begin{itemize}
\item Las diez estrategias acumulan ganancias de forma constante en todo el periodo.
\item No hay interrupciones fuertes en la crisis de 2008, la crisis de deuda europea ni la pandemia.
\item Las líneas grises muestran las demás estrategias como comparación.
\end{itemize}
\vspace{8mm}
{\color{muted}\smallsize Figura tomada de Cakici y Zaremba (2024).}
\end{minipage}
\end{tcolorbox}
\end{minipage}

\vfill
\hspace*{34mm}\begin{minipage}{773mm}
\begin{tcolorbox}[card=teal]\cardtitle[teal]{Resultados}
\begin{tcolorbox}[enhanced, sharp corners, colback=tealSoft, boxrule=0pt, left=14mm, right=14mm, top=10mm, bottom=11mm,
  borderline west={4mm}{0mm}{teal}, fontupper=\bodysize\color{ink}\raggedright]
{\color{teal}\smallsize\bfseries\textls[110]{HALLAZGO CENTRAL}}\par\vspace{4mm}
{\bfseries Los retornos de los países son predecibles, con un R² fuera de muestra de 0.29 \% a 0.95 \%, y pocas variables explican la mayor parte.}
\end{tcolorbox}
\vspace{11mm}
\begin{minipage}[t]{0.48\linewidth}
\begin{itemize}
\item Las variables clave son la reversión de largo plazo, el momentum, el rendimiento de las utilidades y el tamaño del mercado.
\item Las medidas de riesgo país tienen un papel menor.
\end{itemize}
\end{minipage}\hfill
\begin{minipage}[t]{0.48\linewidth}
\begin{itemize}
\item La predictibilidad persiste en mercados pequeños, ilíquidos y poco integrados, y se debilita con el tiempo.
\item La combinación de pronósticos logra 1.46 \% mensual entre el quintil superior y el inferior, con ratio de Sharpe de 1.29.
\end{itemize}
\end{minipage}
\end{tcolorbox}
\end{minipage}

\vfill
\hspace*{34mm}\begin{minipage}{773mm}
\begin{tcbraster}[raster columns=3, raster equal height, raster column skip=16mm]
\begin{tcolorbox}[card, raster multicolumn=2]\cardtitle{Conclusiones}El resultado apoya la idea de segmentación parcial de los mercados. Convertir esa predictibilidad en estrategias rentables es difícil.\end{tcolorbox}
\begin{tcolorbox}[card=gold, colback=goldSoft]\cardtitle{Aporte al tema 61}Confirma que los mercados pequeños e ilíquidos, como la BVL, se comportan de forma distinta a los grandes porque están parcialmente segmentados.\end{tcolorbox}
\end{tcbraster}
\end{minipage}

\vfill
% ======================= PIE =======================
\begin{tcolorbox}[enhanced, sharp corners, colback=navy, boxrule=0pt, width=\paperwidth,
  left=34mm, right=34mm, top=14mm, bottom=11mm,
  overlay={\fill[gold] (frame.north west) rectangle ([yshift=-5mm]frame.north east);}]
\begin{minipage}[c]{245mm}\raggedright
{\color{acento}\smallsize\bfseries\textls[110]{FUENTE DE DATOS DEL ESTUDIO}}\par\vspace{5mm}
{\color{white}\smallsize Datos de 71 mercados bursátiles internacionales, de 1985 a 2021.}
\end{minipage}\hfill
\begin{minipage}[c]{385mm}\raggedright
{\color{acento}\smallsize\bfseries\textls[110]{REFERENCIA}}\par\vspace{5mm}
{\color{white}\smallsize\hangindent=1.2em\hangafter=1 Cakici, N., y Zaremba, A. (2024). What drives stock returns across countries? Insights from machine learning models. \textit{International Review of Financial Analysis}, \textit{96}(Parte A), Artículo 103569. \url{https://doi.org/10.1016/j.irfa.2024.103569}\par}
\end{minipage}\hfill
\begin{minipage}[c]{100mm}\centering
\tcbox[enhanced, sharp corners, colback=white, boxrule=0pt, left=3mm, right=3mm, top=3mm, bottom=3mm]{\qrcode[hyperlink, height=46mm]{https://doi.org/10.1016/j.irfa.2024.103569}}\par\vspace{4mm}
{\color{acento}\fontsize{17}{21}\selectfont\bfseries\textls[100]{ESCANEA Y LEE EL ARTÍCULO}}
\end{minipage}
\end{tcolorbox}
\end{document}
